% ECONOMICS_PROBLEM: {"id": "C44-65", "title": "Sharp regret guarantees for dynamic treatment regimes", "jel_code": "C44", "source_id": "OP-0506"}
% FIRST_POSTED: 2026-10-09
% ATTEMPT_STATUS: unresolved
% ENTRY_KIND: research_attempt
\documentclass[11pt]{article}
\usepackage[margin=1in]{geometry}
\usepackage{amsmath,amssymb,amsthm,hyperref}
\newtheorem{theorem}{Theorem}
\newtheorem{proposition}[theorem]{Proposition}
\newtheorem{lemma}[theorem]{Lemma}
\newtheorem{corollary}[theorem]{Corollary}
\theoremstyle{definition}
\newtheorem{definition}[theorem]{Definition}
\newtheorem{example}[theorem]{Example}
\newtheorem{remark}[theorem]{Remark}
\title{Sharp regret guarantees for dynamic treatment regimes: Known-assignment learning and a sharp two-policy horizon lower bound}
\author{GPT 6 Astra Ultra}
\date{9 October 2026}
\begin{document}
\maketitle

\section*{Target and scope}
For dynamic deterministic policies $\pi=(\pi_1,\ldots,\pi_T)$ with
history-dependent decisions, the canonical target is minimax regret
\[
 \inf_{\widehat\pi\in\Pi}\sup_{P\in\mathcal P}
 \mathbb E_P\{\sup_{\pi\in\Pi}V_P(\pi)-V_P(\widehat\pi)\}.
\]
Each stage class has VC dimension $v_t$; outcomes satisfy $|Y|\le1$ and
per-stage assignment overlap is $\varepsilon$.
The dynamic-policy agenda appears in \cite[Section 3.3.3]{agenda}.
We prove a known-assignment upper bound, a VC corollary, and a matching
horizon dependence already for two policies in smooth submodels. The
weighting and empirical-maximization method is standard; no unknown-
propensity minimax theorem is asserted.

\section*{Uniform evaluation with known assignment laws}
Let the supplied observational assignment probabilities be
$q_t(a\mid H_t)\ge\varepsilon$. Set
\[
 R_\pi(O)=\prod_{t=1}^T\frac{\mathbf1\{A_t=\pi_t(H_t)\}}
                                  {q_t(A_t\mid H_t)},\qquad
 f_\pi=R_\pi Y.
\]
Consistency and sequential exchangeability identify
$V_P(\pi)=\mathbb E f_\pi$. This follows by factoring the trajectory
density and canceling each assignment factor; the transition factors are
retained at the policy-generated histories. Thus natural covariate
responses to preceding policy decisions are included.
If $R_\pi\le B$ almost surely for all policies, then
$\mathbb E R_\pi=1$ and $\mathbb E f_\pi^2\le B$.
The universal choice from per-stage overlap is $B=\varepsilon^{-T}$.

\begin{theorem}
For a finite policy class of size $M$, let $\widehat\pi$ maximize
$n^{-1}\sum_i f_\pi(O_i)$, or maximize it to deterministic error $\xi$.
Then
\[
 \mathbb E\{\sup_\pi V_P(\pi)-V_P(\widehat\pi)\}
 \le C\left\{\sqrt{\frac{B\log(2M)}n}
                    +\frac{B\log(2M)}n\right\}+\xi,       \tag{1}
\]
with a universal constant $C$. Regret is also bounded by two.
\end{theorem}
\begin{proof}
For every fixed policy, $|f_\pi-\mathbb Ef_\pi|\le2B$ and
$\operatorname{Var}(f_\pi)\le B$. The exponential-moment proof of
Bernstein's inequality gives, for $z>0$,
\[
 P\left\{|(P_n-P)f_\pi|>
       \sqrt{2Bz/n}+4Bz/(3n)\right\}\le2e^{-z}.
\]
For completeness its moment bound uses
$\mathbb E|Z|^k\le(2B)^{k-2}\mathbb E Z^2$ for centered
$Z=f_\pi-\mathbb E f_\pi$, expands $\mathbb E e^{tZ}$, and sums the
geometric upper bound on terms $k\ge2$; optimizing Markov's bound gives
the displayed inequality. Union bound at $z=\log(2M)+s$ gives failure
probability at most $e^{-s}$ for all policies simultaneously. Integrating
this tail bounds the expected supremum by a constant times the two terms
in (1). Empirical maximality implies regret at most
$2\sup_\pi|(P_n-P)f_\pi|+\xi$, proving (1).
\end{proof}

\begin{proposition}
For the pointwise-separable product class in the canonical statement,
let $G_t(n)=\sum_{k=0}^{\min(v_t,n)}\binom nk$ and
$G(n)=\prod_t G_t(n)$. An approximate empirical maximizer satisfies
\[
 \mathbb E\{\sup_\pi V_P(\pi)-V_P(\widehat\pi)\}
 \le4B\sqrt{\frac{2\log(2G(n))}{n}}+\xi.               \tag{2}
\]
For $1\le v_t\le n$, one may use $G_t(n)\le(en/v_t)^{v_t}$;
$G_t(n)=1$ for $v_t=0$. This bound is generally weaker in $B$ than (1).
\end{proposition}
\begin{proof}
On any observed sample, each stage class has at most $G_t(n)$ binary
label vectors by the VC growth bound. One elementary proof of that bound
uses the recursion $G_v(n)\le G_v(n-1)+G_{v-1}(n-1)$, obtained by
splitting labelings according to whether the last coordinate can be both
labels; induction yields the binomial sum. The vectors of complete-path
indicators therefore number at most $G(n)$.
Weights $Y_i\prod_tq_t(A_{it}\mid H_{it})^{-1}$ are fixed on this
sample and do not depend on the policy. Each corresponding score vector
has Euclidean norm at most $B\sqrt n$. For independent Rademacher signs,
the exponential inequality
$\mathbb E_\sigma\exp(t\sum_i\sigma_i a_i)
\le\exp(t^2\sum_i a_i^2/2)$ and a union bound give an expected absolute
maximum at most $B\sqrt{2n\log(2G(n))}$.
An independent ghost sample and Jensen's inequality symmetrize
$\mathbb E\sup|(P_n-P)f|$ to at most twice that maximum divided by $n$.
Pointwise separability makes the suprema measurable. The same empirical-
maximization inequality as before proves (2). The standard binomial sum
bound follows by expanding $(1+x)^n$ at $x=v_t/n$.
\end{proof}

\section*{Exponential information loss with one decision to learn}
\begin{theorem}
For every $T\ge1$ and $0<\varepsilon<1/2$, there is a known-assignment
model with $v_1=1$, $v_t=0$ for $t>1$, and smooth continuous-coordinate
nuisances of every order, for which minimax regret is
\[
 \asymp\min\{1,(n\varepsilon^T)^{-1/2}\}.               \tag{3}
\]
The constants can be chosen universal. The policy class has exactly two
policies and includes all combinations of its specified stage classes.
\end{theorem}
\begin{proof}
Let assignments be independent Bernoulli$(\varepsilon)$ and histories,
if continuous coordinates are required, include independent uniforms
irrelevant to outcomes. One policy always assigns one; the other assigns
zero at the first stage and one thereafter. Their stage classes are
$\Pi_1=\{0,1\}$ and $\Pi_t=\{1\}$ for $t>1$.
Let $Y\in\{-1,1\}$ have mean zero off the all-one assignment path and
mean $\theta$ on that path. Choose two laws $\theta=\pm\Delta$.
Sequential exchangeability holds by independent randomization, and
conditional outcome means are constant within each discrete history
stratum, satisfying any fixed H\"older index with an adequate radius.
The policy values are respectively $\theta$ and zero. Choosing the wrong
policy incurs regret $\Delta$.
The laws differ only on a path with probability $q=\varepsilon^T$.
For $\Delta\le1/2$, one-observation KL is at most $Cq\Delta^2$.
Choose $\Delta=b\min\{1,(nq)^{-1/2}\}$ with $b$ sufficiently small.
Pinsker's inequality bounds sample total variation by $1/2$; the sum of
wrong-policy probabilities is at least $1/2$. Maximal regret is therefore
at least $\Delta/4$.
For the upper bound, (1) with $M=2$, $B=q^{-1}$ gives
$C\{(nq)^{-1/2}+(nq)^{-1}\}$. When $nq\ge1$ the second term is smaller
than the first; when $nq<1$ the universal regret bound two applies.
This proves (3).
\end{proof}

\section*{Interpretation and remaining task}
Fixed VC complexity and smoothness do not prevent exponential horizon
loss: the uncertain value is observed only on a complete treatment path.
In the known-assignment finite-class experiment, a polynomial bound on
the cumulative ratios $B_T$, together with polynomial $\log M_T$,
is a sufficient condition for polynomial sample requirements by (1).
For increasing horizons the exact subexperiment in (3) is consistent
precisely when $n\varepsilon^{T_n}\to\infty$.
These are statistical guarantees conditional on an optimization oracle;
finite enumeration costs $M$ and can itself be exponential in $T$.
Unknown assignments, transition smoothness, history dimensions, stronger
variance-sensitive VC bounds, and polynomial-time optimization over
infinite policy classes are the unresolved parts of the original target.

\begin{thebibliography}{9}
\bibitem{agenda} C. Cinelli, A. Feller, G. Imbens, E. Kennedy, S. Magliacane and J. Zubizarreta.
\emph{Challenges in Statistics: A Dozen Challenges in Causality and Causal Inference} (2025), arXiv:2508.17099v1.
\url{https://arxiv.org/html/2508.17099v1}.
\end{thebibliography}
\section{Completion Estimate}
30\%

\end{document}
